% GENERATED FILE. Edit canonical lessons, then run npm run generate:books.
% canonical-source-hash: fnv1a64:5d21867890658dd2
% canonical-lessons: GU-C447-harakh, GU-C447-harkhavu, GU-C447-umalko, GU-C447-anubhavvu, GU-C447-racvu-delight

\chapter{A Morning Full of Joy}
\label{ch:gu-a-morning-full-of-joy}
\hlchaptermodality{447}

\begin{chapteropening}
\textbf{By the end of this chapter:} \emph{I can say how good news makes me feel.}

Complete the last lesson of chapter 447: I can say how good news makes me feel.
\end{chapteropening}

\section[\texorpdfstring{\gu{હરખ} (harakh)}{હરખ (harakh)}]{\gu{હરખ} (harakh) — delight; joy at good news}

\label{lesson:GU-C447-harakh}



\begin{quote}
Before the new one: say the Gujarati for to pester, then the Gujarati for to hope.
\end{quote}

\subsection*{You'll want to know: \gu{હરખ}}
\textbf{\gu{હરખ}} — \emph{harakh} — \textquotedblleft{}delight; joy at good news\textquotedblright{}.

\begin{quote}
\gu{દીકરાનું} \gu{લગ્ન} \gu{નક્કી} \gu{થયું}, \gu{તેથી} \gu{દાદીને} \gu{ઘણો} \gu{હરખ} \gu{થયો}.
\end{quote}

\begin{quote}
\emph{dīkrānũ lagna nakkī thayũ, tethī dādīne ghaṇo harakh thayo.}
\end{quote}

\begin{quote}
My grandmother was filled with delight when her son's wedding was fixed.
\end{quote}

A masculine noun.

\begin{cousinweb}
\textbf{\gu{હરખ}} (\emph{harakh}) probably goes back to Sanskrit \emph{harṣa}, \textquotedblleft{}joy, delight\textquotedblright{}, which Monier-Williams records; it is the everyday spoken shape of that learned word.
\end{cousinweb}

\begin{grammarlens}[title={What you've built}]
One more word for describing what you lived and felt.
\end{grammarlens}

\subsection*{Your turn}
\begin{itemize}
\raggedright
  \item \emph{Say it:} \emph{harakh}
  \item \emph{Say it:} \emph{harakh}, once more
  \item \emph{Say it:} \emph{dīkrānũ lagna nakkī thayũ, tethī dādīne ghaṇo harakh thayo.} — My grandmother was filled with delight when her son's wedding was fixed.
  \item \emph{Recall:} say the Gujarati for to pester, then the Gujarati for to hope, then say \emph{harakh} again
\end{itemize}

\begin{tcolorbox}[breakable,colback=teal!4,colframe=teal!35!black,title={Before you move on}]
What does \gu{હરખ} mean? (\textquotedblleft{}Delight; joy at good news\textquotedblright{}.) Say it once more.
\end{tcolorbox}
\section[\texorpdfstring{\gu{હરખાવું} (harkhāvũ)}{હરખાવું (harkhāvũ)}]{\gu{હરખાવું} (harkhāvũ) — to rejoice; to be delighted}

\label{lesson:GU-C447-harkhavu}



\begin{quote}
Before the new one: say the Gujarati for to hope, then the Gujarati for delight.
\end{quote}

\subsection*{You'll want to know: \gu{હરખાવું}}
\textbf{\gu{હરખાવું}} — \emph{harkhāvũ} — \textquotedblleft{}to rejoice; to be delighted\textquotedblright{}.

\begin{quote}
\gu{પરિણામ} \gu{સાંભળીને} \gu{મીરા} \gu{હરખાઈ} \gu{ગઈ}.
\end{quote}

\begin{quote}
\emph{pariṇām sā̃bhḷīne Mīrā harkhāī gaī.}
\end{quote}

\begin{quote}
Mira was delighted when she heard the result.
\end{quote}

Past (she …): \gu{તે} \gu{હરખાઈ} \gu{ગઈ}, \emph{te harkhāī gaī}.

\begin{cousinweb}
\textbf{\gu{હરખાવું}} (\emph{harkhāvũ}) is built on \textbf{\gu{હરખ}} (\emph{harakh}), \textquotedblleft{}delight\textquotedblright{}, with the verb ending \emph{-āvũ}; the noun probably goes back to Sanskrit \emph{harṣa}, \textquotedblleft{}joy\textquotedblright{} (Monier-Williams).
\end{cousinweb}

\begin{grammarlens}[title={What you've built}]
One more word for describing what you lived and felt.
\end{grammarlens}

\subsection*{Your turn}
\begin{itemize}
\raggedright
  \item \emph{Say it:} \emph{harkhāvũ}
  \item \emph{Say it:} \emph{harkhāvũ}, once more
  \item \emph{Say it:} \emph{pariṇām sā̃bhḷīne Mīrā harkhāī gaī.} — Mira was delighted when she heard the result.
  \item \emph{Recall:} say the Gujarati for to hope, then the Gujarati for delight, then say \emph{harkhāvũ} again
\end{itemize}

\begin{tcolorbox}[breakable,colback=teal!4,colframe=teal!35!black,title={Before you move on}]
What does \gu{હરખાવું} mean? (\textquotedblleft{}To rejoice; to be delighted\textquotedblright{}.) Say it once more.
\end{tcolorbox}
\section[\texorpdfstring{\gu{ઉમળકો} (umaḷko)}{ઉમળકો (umaḷko)}]{\gu{ઉમળકો} (umaḷko) — warmth of feeling; eager affection}

\label{lesson:GU-C447-umalko}



\begin{quote}
Before the new one: say the Gujarati for delight, then the Gujarati for to rejoice.
\end{quote}

\subsection*{You'll want to know: \gu{ઉમળકો}}
\textbf{\gu{ઉમળકો}} — \emph{umaḷko} — \textquotedblleft{}warmth of feeling; eager affection\textquotedblright{}.

\begin{quote}
\gu{માસીએ} \gu{ઉમળકાથી} \gu{બારણું} \gu{ખોલ્યું}.
\end{quote}

\begin{quote}
\emph{māsīe umaḷkāthī bārṇũ kholyũ.}
\end{quote}

\begin{quote}
My aunt opened the door with warm delight.
\end{quote}

A masculine noun.

\begin{cousinweb}
\textbf{\gu{ઉમળકો}} (\emph{umaḷko}) goes back to a reconstructed \emph{*ummal-}, \textquotedblleft{}rise, bubble up\textquotedblright{}: a feeling that wells over. CDIAL lists the Gujarati word as \textquotedblleft{}emotion, ardent love\textquotedblright{} (CDIAL 2346).
\end{cousinweb}

\begin{grammarlens}[title={What you've built}]
One more word for describing what you lived and felt.
\end{grammarlens}

\subsection*{Your turn}
\begin{itemize}
\raggedright
  \item \emph{Say it:} \emph{umaḷko}
  \item \emph{Say it:} \emph{umaḷko}, once more
  \item \emph{Say it:} \emph{māsīe umaḷkāthī bārṇũ kholyũ.} — My aunt opened the door with warm delight.
  \item \emph{Recall:} say the Gujarati for delight, then the Gujarati for to rejoice, then say \emph{umaḷko} again
\end{itemize}

\begin{tcolorbox}[breakable,colback=teal!4,colframe=teal!35!black,title={Before you move on}]
What does \gu{ઉમળકો} mean? (\textquotedblleft{}Warmth of feeling; eager affection\textquotedblright{}.) Say it once more.
\end{tcolorbox}
\section[\texorpdfstring{\gu{અનુભવવું} (anubhavvũ)}{અનુભવવું (anubhavvũ)}]{\gu{અનુભવવું} (anubhavvũ) — to feel; to experience}

\label{lesson:GU-C447-anubhavvu}



\begin{quote}
Before the new one: say the Gujarati for to rejoice, then the Gujarati for warmth of feeling.
\end{quote}

\subsection*{You'll want to know: \gu{અનુભવવું}}
\textbf{\gu{અનુભવવું}} — \emph{anubhavvũ} — \textquotedblleft{}to feel; to experience\textquotedblright{}.

\begin{quote}
\gu{લગ્નમાં} \gu{મેં} \gu{ઘણી} \gu{ખુશી} \gu{અનુભવી}.
\end{quote}

\begin{quote}
\emph{lagnamā̃ mẽ ghaṇī khuśī anubhavī.}
\end{quote}

\begin{quote}
I felt great happiness at the wedding.
\end{quote}

Past, with \gu{તેણે} (he or she …): \gu{તેણે} \gu{ખુશી} \gu{અનુભવી}, \emph{teṇe khuśī anubhavī}.

\begin{cousinweb}
\textbf{\gu{અનુભવવું}} (\emph{anubhavvũ}) is built on \textbf{\gu{અનુભવ}} (\emph{anubhav}), \textquotedblleft{}experience\textquotedblright{}, a tatsama of Sanskrit \emph{anubhava}, \textquotedblleft{}perception, fruition\textquotedblright{} (Monier-Williams), with the verb ending \emph{-vũ}.
\end{cousinweb}

\begin{grammarlens}[title={What you've built}]
One more word for describing what you lived and felt.
\end{grammarlens}

\subsection*{Your turn}
\begin{itemize}
\raggedright
  \item \emph{Say it:} \emph{anubhavvũ}
  \item \emph{Say it:} \emph{anubhavvũ}, once more
  \item \emph{Say it:} \emph{lagnamā̃ mẽ ghaṇī khuśī anubhavī.} — I felt great happiness at the wedding.
  \item \emph{Recall:} say the Gujarati for to rejoice, then the Gujarati for warmth of feeling, then say \emph{anubhavvũ} again
\end{itemize}

\begin{tcolorbox}[breakable,colback=teal!4,colframe=teal!35!black,title={Before you move on}]
What does \gu{અનુભવવું} mean? (\textquotedblleft{}To feel; to experience\textquotedblright{}.) Say it once more.
\end{tcolorbox}
\section[\texorpdfstring{\gu{રાચવું} (rācvũ)}{રાચવું (rācvũ)}]{\gu{રાચવું} (rācvũ) — to delight in; to revel in}

\label{lesson:GU-C447-racvu-delight}



\begin{quote}
Before the new one: say the Gujarati for warmth of feeling, then the Gujarati for to experience.
\end{quote}

\subsection*{You'll want to know: \gu{રાચવું}}
\textbf{\gu{રાચવું}} — \emph{rācvũ} — \textquotedblleft{}to delight in; to revel in\textquotedblright{}.

\begin{quote}
\gu{બાળકો} \gu{નવાં} \gu{રમકડાંમાં} \gu{રાચે} \gu{છે}.
\end{quote}

\begin{quote}
\emph{bāḷko navā̃ ramakḍā̃mā̃ rāce che.}
\end{quote}

\begin{quote}
The children are delighting in their new toys.
\end{quote}

Past (he …): \gu{તે} \gu{રમતમાં} \gu{રાચ્યો}, \emph{te ramatmā̃ rācyo}.

\begin{cousinweb}
\textbf{\gu{રાચવું}} (\emph{rācvũ}) is inherited from Sanskrit \emph{rajyatē}, \textquotedblleft{}takes pleasure in\textquotedblright{}; CDIAL lists the Gujarati verb, \textquotedblleft{}to be pleased\textquotedblright{} (CDIAL 10584).
\end{cousinweb}

\begin{grammarlens}[title={What you've built}]
One more word for describing what you lived and felt.
\end{grammarlens}

\subsection*{Your turn}
\begin{itemize}
\raggedright
  \item \emph{Say it:} \emph{rācvũ}
  \item \emph{Say it:} \emph{rācvũ}, once more
  \item \emph{Say it:} \emph{bāḷko navā̃ ramakḍā̃mā̃ rāce che.} — The children are delighting in their new toys.
  \item \emph{Recall:} say the Gujarati for warmth of feeling, then the Gujarati for to experience, then say \emph{rācvũ} again
\end{itemize}

\begin{tcolorbox}[breakable,colback=teal!4,colframe=teal!35!black,title={Before you move on}]
What does \gu{રાચવું} mean? (\textquotedblleft{}To delight in; to revel in\textquotedblright{}.) Say it once more.
\end{tcolorbox}
